% docs/manual/operations/appendices/D-performance-guidance.tex
% 附录 D：性能建议（建议性指导）

\chapter{性能建议（建议性指导）}

本附录给出 \openbench{} 在不同负载下的资源画像与调优建议。\textbf{不是实测基准}（spec §2 决定的方向）；具体数字会因机器、数据集、模型差异较大。请把它当作初次部署时的"起点"，运维过程中自己测自己机器。

\section{资源画像（按负载分类）}

\subsection{Tiny：单变量月评估}

\begin{itemize}
  \item 配置：1 变量、1 sim、1 ref、5 年、月尺度、全球 0.5°
  \item 时间：5-15 分钟（笔记本 4 核）
  \item 内存：2-6 GB 峰值
  \item 磁盘：< 500 MB
\end{itemize}

\textbf{适用机器}：笔记本、入门工作站。

\subsection{Small：典型科研评估}

\begin{itemize}
  \item 配置：4 变量、2 sim、2 ref、5 年、月尺度、全球 0.5°
  \item 时间：30 分钟 - 2 小时（工作站 16 核）
  \item 内存：8-20 GB 峰值
  \item 磁盘：1-5 GB
\end{itemize}

\textbf{适用机器}：工作站。

\subsection{Medium：多模型对比}

\begin{itemize}
  \item 配置：8 变量、5 sim、5 ref、10 年、月尺度、0.25°
  \item 时间：3-12 小时（HPC 32 核）
  \item 内存：30-80 GB 峰值
  \item 磁盘：10-30 GB
\end{itemize}

\textbf{适用机器}：HPC 节点（中等规模）。

\subsection{Large：千站点 + 高频}

\begin{itemize}
  \item 配置：6 变量、3 sim、~500 站、5 年、小时尺度
  \item 时间：12-48 小时（HPC 64 核）
  \item 内存：50-150 GB 峰值
  \item 磁盘：30-100 GB
\end{itemize}

\textbf{适用机器}：HPC 大节点。

\subsection{XL：全球高频长时段}

\begin{itemize}
  \item 配置：10+ 变量、多 sim、全球 daily 或 hourly、10 年
  \item 时间：1-7 天（HPC 128 核）
  \item 内存：100-500 GB 峰值
  \item 磁盘：100-500 GB
\end{itemize}

\textbf{适用机器}：HPC 大内存节点。考虑分批跑（按变量分组多次）。

\section{粗略时间公式}

时间 ≈ \(\frac{V \times S \times R \times Y \times G}{\alpha \times C}\)

\begin{itemize}
  \item \(V\)：变量数
  \item \(S\)：simulation 数
  \item \(R\)：reference 数（每变量平均）
  \item \(Y\)：年数
  \item \(G\)：分辨率因子（hourly=24, daily=24, monthly=1）
  \item \(C\)：核数
  \item \(\alpha\)：机器系数（笔记本=1, 工作站=2, HPC=4）
\end{itemize}

不严谨。实测后调 \(\alpha\)。

\section{典型错误：HPC 上 num\_cores=1}

最常见的"性能浪费"：HPC 分配了 32 核但 \openbench{} 只用 1。原因：

\begin{itemize}
  \item 没设 \cli{--cores} 或 yaml \yamlkey{num\_cores}
  \item 默认 \texttt{os.cpu\_count()} 在 SLURM 节点上看到 \emph{物理核} 而不是 \emph{分配核}（取决于 cgroup 配置）
\end{itemize}

\textbf{修法}：永远显式设：

\begin{minted}{bash}
openbench run config.yaml --cores $SLURM_CPUS_PER_TASK
\end{minted}

\section{调参第一性原理}

不要凭直觉调；按顺序：

\begin{enumerate}
  \item \textbf{Profile}：先用 \texttt{/usr/bin/time -v} 跑一次 small 配置，看 CPU\% / 内存峰值 / wall time
  \item \textbf{识别瓶颈}：CPU bound？内存？IO？
  \item \textbf{对症调}：见第 4 章
  \item \textbf{再 profile}：调完再跑同样负载，看变化
  \item \textbf{Scale up}：把改动应用到大评估
\end{enumerate}

\section{该做与不该做}

\textbf{该做}：

\begin{itemize}
  \item 先 small 后 large，验证流程
  \item 显式 num\_cores，不靠默认
  \item Reference 数据放高速盘
  \item 启用增量 cache
  \item HPC 上用 SLURM 不直接 login 节点跑长任务
  \item 设 OMP/MKL\_NUM\_THREADS=1
\end{itemize}

\textbf{不该做}：

\begin{itemize}
  \item HPC 上启动 GUI
  \item 一次性跑全部变量（先 1 变量 + 1 sim 验证）
  \item 直接在 \texttt{\$HOME} 跑（共享 NFS 慢）
  \item 大评估不开监控
  \item OOM 后无脑加 num\_cores（错的方向，应该减）
\end{itemize}

\section{何时该升级硬件}

如果按上述调优后，运行时间仍：

\begin{itemize}
  \item > 3 天单次评估 → 加节点 / 加内存
  \item 内存反复 OOM → 升级到大内存节点（\textgreater 256 GB）
  \item NFS 抖动严重 → 把数据搬到 Lustre
\end{itemize}

\section{未来方向}

\openbench{} v3.0a1 的并行还有提升空间：

\begin{itemize}
  \item Task 级（变量级）并发尚未启用
  \item Comparison/statistics 阶段串行
  \item Visualization 串行
\end{itemize}

后续版本启用后，相同硬件预期吞吐能提升 2-5x。

\warninline{本附录是\textbf{建议性}，非实测。具体到你的机器与负载，请自己 profile 一次再决定参数。}
